%======================================================================
\section{Introduction}\label{sec:intro}
%======================================================================

Additional Tier~1 (AT1) bonds, often called contingent convertibles or CoCos, are the
most equity-like instruments that European banks issue as debt. They count as
going-concern Tier~1 capital under the Capital Requirements Regulation, they are
perpetual and callable only by the issuer, their coupons are fully discretionary and
non-cumulative, and their principal is converted into shares or written down when the
bank's Common Equity Tier~1 (CET1) ratio falls below a contractual trigger. In
addition, the resolution authority can impose losses on them at the point of
non-viability (PONV), whatever the level of the trigger. The write-down of
CHF~16 billion of Credit Suisse AT1s in March 2023, while the bank reported a CET1
ratio around $14\%$, made the last feature the dominant one in market pricing.

Every one of these cash-flow contingencies is written, directly or indirectly, on the
bank's capital. The trigger is a CET1 level; the coupon restriction of the Maximum
Distributable Amount (MDA) regime is triggered when CET1 falls inside the combined
buffer requirement; regulatory approval of a call requires the bank to remain above its
requirements after redemption; and the regulator's decision to declare the bank
failing or likely to fail is, in practice, informed by the same deterioration of
capital and market access. This paper takes that observation literally: it models the
AT1 bond as a claim on the CET1 ratio, rather than as a straight bond with an equity
barrier attached.

\paragraph{Why not an equity barrier.}
The workhorse market model for CoCos treats conversion as a barrier event on the share
price \citep{DeSpiegeleerSchoutens2012}. That approach is convenient, because the share
price is observed and its volatility is quoted, but it needs a mapping from the share
price to the accounting ratio on which the contract is written. That mapping is not
observable and not stable: banks' equity traded at a fraction of book value in 2023
with unchanged capital ratios, and at a premium in other periods. A model written on
the CET1 ratio makes the trigger, the MDA restriction and the call constraint
consistent by construction. Its cost is that CET1 is not traded and is observed
quarterly with a lag, so its dynamics must be estimated under the physical measure and
the residual risk premia must be absorbed elsewhere. We absorb them in a write-down
hazard that is implied from bond prices, and we then ask how much of that hazard the
bank's other bonds also price.

\paragraph{Capital alone is not enough.}
A model in which CET1 is the only state variable cannot price BNP Paribas's AT1s. With
the historical volatility of BNP's CET1 ratio, the trigger is $7.8$ percentage points
away and is almost never reached in simulation: the bonds would be worth $114$--$118$
against market prices of $94$--$98$. The call decision, moreover, depends on the cost
of refinancing in the market, which moves with sector spreads and not only with the
bank's capital. We therefore add a second state variable: the issuer's new-issue AT1
spread, a mean-reverting Gaussian process correlated with capital shocks, with a
capital add-on once CET1 enters the buffer. This spread does two jobs. It drives the
call decision, so that the issuer extends exactly when refinancing is expensive, and it
scales a doubly-stochastic write-down hazard, so that extension states are also the
states in which the bond is riskier. Without the second link, extension would be valued
with a discount rate that does not reflect the state, and the option not to call could
even look valuable to the investor; we show this explicitly in
Section~\ref{sec:robust}.

\paragraph{What the hazard prices.}
The implied write-down hazard absorbs everything the mechanical features leave
unexplained, and on its own it cannot tell the bank's credit risk apart from risk that is
specific to the instrument. The bank's other bonds can. Under the EU resolution
framework the creditor hierarchy is respected: CET1 absorbs losses first, then AT1, then
Tier~2. The joint statement of the Single Resolution Board, the EBA and ECB Banking
Supervision of 20 March 2023 reaffirmed this ordering after the Credit Suisse case.
Tier~2 bonds are therefore exposed to the same resolution event as AT1, at a better
recovery, but have none of the AT1's coupon discretion, MDA restriction, trigger or
perpetual extension. Pricing them on the same simulated state identifies how much of the
AT1 hazard is compensation that Tier~2 also earns.

\subsection{Contributions}

The paper makes five contributions.
\begin{enumerate}
  \item \textbf{A capital-ratio model with an endogenous call and a state-dependent
  hazard.} The CET1 ratio follows a mean-reverting jump-diffusion around the bank's
  target, the refinancing spread a correlated Ornstein--Uhlenbeck process, and the
  write-down time is a Cox process whose intensity is proportional to the effective
  spread. The MDA threshold is time-varying, which matters because scheduled buffer
  step-ups shrink the bank's cushion at constant capital.
  \item \textbf{An exact conditional Monte Carlo for the hazard.} Conditioning on the
  state paths integrates the write-down time out analytically. The price becomes an
  explicit function $V(\lambda)$ that is smooth, strictly decreasing and strictly convex,
  so the implied hazard is unique and is found by a bracketing root-finder on a single
  set of paths. The same construction prices Tier~2 bonds with a recovery leg and a
  piecewise term structure of the hazard.
  \item \textbf{A Tier~2 anchor for the AT1 hazard.} Writing the AT1 hazard as
  $\lambda=\lres+\lspec$, with two independent exponential clocks, gives the AT1 and the
  Tier~2 the right joint survival structure. Calibrating $\lres$ to the bank's Tier~2 and
  checking it against the senior non-preferred bond turns an unexplained premium into an
  observable part and an AT1-specific residual.
  \item \textbf{An order-robust attribution of the spread.} Sequential decompositions of
  a credit spread into its components depend on the order in which features are switched
  on. We compute the exact Shapley value over all orders and a Shapley value within
  blocks, and show which conclusions survive. The order matters for the extension
  component and not for the split between Tier~2-priced and AT1-specific risk.
  \item \textbf{An application with honest identification.} We calibrate the model to
  BNP Paribas's three USD AT1s, two Tier~2 bonds, one senior non-preferred bond and its
  senior CDS, with 50 quarters of CET1 history. We report Monte Carlo errors, price the
  identification limits of evaluated bid--offer quotes, bootstrap the capital dynamics
  and show which results depend on the weakly identified speed of capital mean
  reversion.
\end{enumerate}

\subsection{Main findings}

On 25 September 2026 the three AT1s trade at yield-to-call spreads of $269$--$280$bp
over call-matched Treasuries. The model reprices them with hazards of $2.21\%$,
$2.18\%$ and $2.19\%$ a year. The dispersion of $2.6$bp of hazard is about four
Monte Carlo standard errors but only $0.20$--$0.22$ points of price, the size of the
bid--offer: the bonds are consistent with a single hazard within the precision of the
data, and cannot discriminate between model variants that differ by a few basis points.

BNP's Tier~2 2035NC34 implies a resolution hazard of $1.72\%$ at the conventional
$20\%$ recovery, and the senior non-preferred bond $1.63\%$ at $40\%$: the two
instruments that share the resolution event agree under the Markit recovery
conventions, and the Tier~2 recovery that equalises them is $16\%$. The AT1 hazard
therefore splits into $1.72\%$ priced in Tier~2 and $0.47\%$ specific to the AT1. In
spread terms, about $180$bp of the AT1 spread is priced in Tier~2, $50$bp is an
AT1-specific premium, and $40$--$46$bp pays for the contract features: expected MDA
coupon cuts ($\approx30$bp, an upper bound), mechanical conversion ($\approx4$bp)
and extension ($6$--$15$bp, decreasing in the reset margin). Because the MDA component
is an upper bound and the total is fixed by the market price, the AT1-specific premium
is a lower bound.

The split survives three robustness checks. A reputational call rule raises the
probability of a call at the first reset from $49$--$58\%$ to $77$--$81\%$, close to
the levels practitioners quote, but moves neither the hazard nor its split. Every
attribution that prices the hazard before the contract features gives the same
$181/50$bp split. A two-piece resolution hazard bootstrapped from the two Tier~2 bonds
($1.40\%$ to 2030, $2.10\%$ after) changes the Tier~2-priced component by less than
5bp.

Capital risk is a cliff, not a slope. At BNP's $13\%$ CET1 ratio the price moves by
about $0.6$ points per point of capital; at the MDA threshold it moves by about four.
Crossing the threshold switches on three mechanisms at once: coupons are cut, the call
is blocked and the hazard rises. The distance that matters for risk management is the
distance to the MDA threshold, which the scheduled O-SII step-ups reduce from $2.5$ to
$2.0$ percentage points by 2028 at constant capital, not the distance to the trigger.
The level of the capital stress depends strongly on the speed of capital mean
reversion, which fifty quarters of history identify only to within a factor of eight;
across that interval the price impact of the EBA adverse scenario ranges from $-7$ to
$-38$ points.

\subsection{Related literature}

\paragraph{Contingent capital.} The academic literature on contingent capital starts
from structural models in which the trigger is written on the value of the bank's
assets or on an accounting ratio. \citet{Pennacchi2010} models bank assets as a
jump-diffusion and prices CoCos with a trigger on the ratio of equity to assets;
\citet{GlassermanNouri2012} derive closed-form prices for contingent capital with a
capital-ratio trigger and gradual conversion. \citet{Flannery2014} surveys the design
issues, including the choice between market and accounting triggers and the role of
regulatory discretion. These models are built to study incentives and design, and
typically do not include the MDA restriction, the issuer's call option or a regulatory
write-down hazard independent of the trigger. Market practice instead relies on
reduced-form or equity-barrier models \citep{DeSpiegeleerSchoutens2012}, in which the
trigger is approximated by a barrier on the share price. \citet{AvdjievEtAl2013}
describe the market and the instrument features in detail. Our model sits between the
two: the state is the accounting ratio on which the contract is written, as in the
structural tradition, but the regulatory and market features that dominate AT1 pricing
today (discretionary coupons, MDA, issuer calls, PONV) are modelled explicitly and the
residual risk is priced in reduced form.

\paragraph{Reduced-form credit risk.} The write-down time is a doubly-stochastic
(Cox) process in the sense of \citet{Lando1998}, with intensity driven by the state, as
in the reduced-form models of \citet{JarrowTurnbull1995} and \citet{DuffieSingleton1999};
\citet{BieleckiRutkowski2002} give the general construction. Our use of the
construction is specific in two respects. First, the intensity is proportional to the
issuer's simulated refinancing spread, which also drives the call, so that the hazard
and the extension decision are consistent state by state. Second, we split the
intensity into two independent clocks to price the AT1 and the Tier~2 jointly, which
is the reduced-form counterpart of the creditor hierarchy.

\paragraph{Monte Carlo methods.} Integrating the default time out conditionally on the
state paths is a form of conditional Monte Carlo, or Rao--Blackwellisation
\citep{Glasserman2004}; it removes the variance due to the default event and makes the
price smooth in the hazard. Discrete monitoring of the trigger introduces the usual
bias of barrier options, which \citet{BroadieGlassermanKou1997} correct with a shift of
the barrier; we show that the correction is negligible here.

\paragraph{Decomposition and attribution.} Attributing a spread to its components by
switching features on one at a time is order-dependent whenever the features interact.
The Shapley value \citep{Shapley1953} resolves the dependence by averaging over all
orders, and \citet{Shorrocks2013} applies it to decompositions in economics; the value
for games with a priori unions \citep{Owen1977} handles groups of related factors. We
use both ideas and show that, for AT1 bonds, the economically consistent attribution
prices the hazard before the contract features.

\subsection{Outline}

Section~\ref{sec:institutions} describes the instrument, the regulatory framework and
the three BNP Paribas bonds. Section~\ref{sec:model} sets out the model and derives
the properties of the state processes. Section~\ref{sec:valuation} derives the pricing
equations, the conditional Monte Carlo estimator and its properties, and the pricing of
Tier~2 and senior bonds. Section~\ref{sec:attribution} defines the spread attribution.
Section~\ref{sec:calibration} describes the data and the calibration.
Section~\ref{sec:results} reports the results, Section~\ref{sec:robust} the robustness
checks and Section~\ref{sec:discussion} discusses implications and limitations. The
appendices contain proofs, numerical diagnostics and the parameter table.
